\section{State probabilities}
\label{sec:state-probs}

Using formula (\ref{stateFormula}), we can obtain the individual state
probabilities. The first, $p_0(t)$, differs between the catastrophe and
killing cases:
$$G(0 , t) = \begin{cases}
\frac{ab(1-\ee^{\beta(a-b) t})  } {  (b-a\ee^{\beta(a-b) t})   }, & \text{(catastrophe)}\\
1 - I(t) + \frac{ab(1-\ee^{\beta(a-b) t})  } {  (b-a\ee^{\beta(a-b) t})   }, & \text{(killing).}
\end{cases}$$
As stated above, the fraction term in each case represents the probability
of reaching $0$ by the death of every individual, and the $1-I(t)$ pertains
to reaching $0$ by collapse in the killing convention. \par
From here, the state probabilities $p_1(t)$, $p_2(t)$, $p_3(t)\dots$ do not
differ between the cases. Repeated differentiation of $G(z,t)$ with respect
to $z$ reveals the formula,
$$\frac{\partial^n G}{\partial z^n} =n! \ee^{\beta(a-b)t}(a-b)^2  \frac{\lrb{1-\ee^{\beta(a-b)t}}^ {n-1} }{\lrb{b - a\ee^{\beta(a-b)t} - \lrb{1-\ee^{\beta(a-b)t}}z} ^ {n+1}} $$
Hence, setting $z=0$, and dividing by $n!$, the state probabilities for
$n\geq1$ are described by,
\begin{equation}
\label{stateprob}
p_n(t) =\lrb{ \frac{ a-b }{{b - a\ee^{\beta(a-b)t}}}}^2 \ee^{\beta(a-b)t} \lrb{\frac{{1-\ee^{\beta(a-b)t}}}{{b - a\ee^{\beta(a-b)t}} } }^{n-1}.
\end{equation}

\subsection{Geometric at every time}
\label{sec:geometric-every-time}

There is more structure here than (\ref{stateprob}) first suggests. Writing
\begin{equation}
P(t):=\frac{1-\ee^{\beta(a-b)t}}{b-a\ee^{\beta(a-b)t}},
\label{eq:Pdef}
\end{equation}
the entire $n$-dependence of $p_n(t)$ lives in the single factor
$P(t)^{n-1}$:
\begin{equation}
p_n(t)=p_1(t)\,P(t)^{n-1},\qquad n\ge1.
\end{equation}
That is, \emph{at every time $t$}, conditional on the population being
non-zero, the load is geometrically distributed with ratio $P(t)$. The ratio
itself evolves monotonically, with
\begin{equation}
P(0)=0,\qquad P(t)\xrightarrow[t\to\infty]{}\frac{1}{a},
\end{equation}
so the transient geometric laws slide towards a fixed limit. This limit is
the heart of \S\ref{sec:qsd}--\S\ref{sec:burst-time-size}.

\subsection{Normalisation: the sum is $\Ihat$}
\label{sec:sum-ihat}

A useful consistency check is to sum \eqref{stateprob} over the positive
states. As a geometric series,
$$
\sum_{n=1}^{\infty}p_n(t)
=\frac{p_1(t)}{1-P(t)}.
$$
Now
$$
1-P(t)=\frac{b-a\ee^{\beta(a-b)t}-1+\ee^{\beta(a-b)t}}{b-a\ee^{\beta(a-b)t}}
=\frac{-(1-b)-\ee^{\beta(a-b)t}(a-1)}{b-a\ee^{\beta(a-b)t}}
=\frac{B+A\ee^{\beta(a-b)t}}{a\ee^{\beta(a-b)t}-b},
$$
and therefore
\begin{equation}
\sum_{n=1}^{\infty}p_n(t)
=\frac{(a-b)^2\ee^{\beta(a-b)t}}
{\lrb{B+A\ee^{\beta(a-b)t}}\lrb{a\ee^{\beta(a-b)t}-b}}
=\Ihat(t).
\label{eq:pnsum}
\end{equation}
The sum of the positive state probabilities is exactly the corrected
non-fixation probability of \S2.4 --- the probability of being neither
ruptured nor extinct. Equation \eqref{eq:pnsum} thereby confirms, by a
completely independent route, the correction of Remark~\ref{rem:ihat}: the
state probabilities of this chapter were already implying the right
$\Ihat(t)$.
